\paragraph{2.1. асимптотические обозначения}
\subparagraph{3.1.1}
Пусть $f(n)$ и $g(n)$ — асимптотически неотрицательные функции. Докажите с помощью базового определения $\Theta$-обозначений, что
\[
\max(f(n), g(n)) = \Theta(f(n) + g(n)).
\]

\subparagraph{3.1.2}
Покажите, что для любых действительных констант $a$ и $b$, где $b > 0$, выполняется соотношение
\[
(n + a)^b = \Theta(n^b).
\]

\subparagraph{3.1.3}
Поясните, почему утверждение «время работы алгоритма $A$ равно как минимум $O(n^2)$» лишено смысла.

\subparagraph{3.1.4}
Справедливы ли соотношения
\[
2^{n+1} = O(2^n)
\]
и
\[
2^{2n} = O(2^n)?
\]

\subparagraph{3.1.5}
Докажите теорему 3.1.

\subparagraph{3.1.6}
Докажите, что время работы алгоритма равно $\Theta(g(n))$ тогда и только тогда, когда его время работы в наилучшем случае равно $O(g(n))$, а в наихудшем — $\Omega(g(n))$.

\subparagraph{3.1.7}
Докажите, что множество $o(g(n)) \cap \omega(g(n))$ является пустым.

\subparagraph{3.1.8}
Можно обобщить наши обозначения на случай двух параметров $n$ и $m$, которые могут возрастать до бесконечности по отдельности с разными скоростями. Для данной функции $g(n, m)$ обозначим как $O(g(n, m))$ множество функций
\[
O(g(n, m)) = \{ f(n, m) : \text{существуют положительные константы } c, n_0 \text{ и } m_0, \text{ такие, что } 0 \leq f(n, m) \leq c g(n, m) \}
\]
для всех $n \geq n_0$ или $m \geq m_0$.